% ============================================================================
%  Muestra del formato. Tabla del Artículo 46 con tres filas copiadas tal cual
%  de Informe/Tabla-Art46-Informe2.md, solo para probar el componente. Se
%  separó el número de la observación en su propia columna y la negrita de
%  Markdown pasó a \textbf. El texto no se tocó.
% ============================================================================
\begin{tablaObservaciones}{Informe 1}
  \grupoObservaciones{General, formalidad y cumplimiento de instrucciones}
  \observacion[sub=0]{01}{Formulario T-12 ausente pese a que el Subdocumento 3 p.6 dice que acompaña la oferta}{Se acepta. Se adjunta el catálogo completo}{S3, anexo Formulario T-12}
  \grupoObservaciones{Subdocumento 1, presentación de la empresa}
  \observacion[sub=1]{17}{Cita normativa errada, el DS 43 del Ministerio de Salud se invoca como norma de transporte}{\textbf{Se aclara.} El texto de la p.8 dice «transporte y almacenamiento de sustancias peligrosas (Decreto Supremo N.º 298 ... y Decreto Supremo N.º 43 ...)». La enumeración es copulativa: el 298 rige el transporte y el 43 el almacenamiento en los patios químicos y estaciones de carga donde inician y terminan los ciclos de las 18 unidades de sustancias peligrosas. La bibliografía identifica al 43 como reglamento de almacenamiento. Se reformula el párrafo para que la distinción quede explícita}{S1 1.2}
  \grupoObservaciones{Subdocumento 2, comprensión del problema y de la necesidad}
  \observacion[sub=2]{28}{No hay un solo dato investigado fuera del caso, la bibliografía los lista y el texto no los usa}{\textbf{Se aclara.} El cuerpo del subdocumento utiliza la Ley N.º 20.123 sobre subcontratación seis veces, el Artículo 25 bis cinco, la Ley de Tránsito N.º 18.290 cinco, el tacógrafo digital seis y el Marco GLEC cinco con su fuente. La Ley N.º 20.123 no aparece en ninguno de los tres documentos de la licitación, de modo que es aporte externo del proponente. Se acepta que el uso no siempre deriva en una consecuencia de diseño y se corrige}{S2 1.1 y 1.2}
\end{tablaObservaciones}
